
\subsubsection{Inference: One-shot Generation}
During training, each conditioning input (video, audio context, text) is dropped-out independently with some probabilities.
This enables the model to perform
\begin{enumerate*}[label=(\arabic*)]
    \item video-to-audio (V2A) generation (dropping out text and audio context),
    \item text-instructed video-to-audio (TV2A) generation (dropping out audio context),
    \item video-to-audio infilling or extension (dropping out text), and
    \item text-instructed video-to-audio infilling or extension,
\end{enumerate*}
with a single model by simply changing the conditioned inputs.

\subsubsection{Inference: Audio Extension}\label{sec:audio_ext}
Due to memory constraint and training efficiency considerations, training data is capped at a predetermined length.
To generate high quality and coherent long-form audio for videos whose lengths are beyond the cap,
we consider two algorithms: \textit{segment-level autoregressive generation} and \textit{multi-diffusion}~\citep{bar2023multidiffusion}.

Given a long video, we first split the video into \textit{overlapping} segments, and assume text captions are available for all segments.
At a high level, both algorithms run inference on individual segments first, and then consolidate the prediction.
Information can propagate across segments through overlapping frames.
Without loss of generality, we assume each segment is $n_{win}$ frames long,
and the end times of consecutive segments differ by $n_{hop}$ frames.
Two consecutive segments overlap by $n_{ctx} = n_{win} - n_{hop}$ frames.
For a video $\bc_{vid}$ of $N$ frames,
there will be $J = \lceil N / n_{hop} \rceil$ segments,
and the $j$-th segment spans $[n_{start}^{(j)}, n_{end}^{(j)}] = [max(0, (j-1)n_{hop} - n_{ctx}), min(N, jn_{hop})]$ where $j \in [J]$.
We denote the text caption for $j$-th segment as $\bc_{txt}^{(j)}$,
and the consolidated prediction at flow step $t_i$ as $\bx_{t_i}^{(j)}$,
assuming the same $T$-step flow time schedule $\{t_i\}_{i=1}^T$ are used for all segments ($t_1 = 0$, $t_T = 1$, and $t_i < t_{i+1}\;\forall i$).
Note that $\bx_0^{(j)}$ is the noise drawn from the prior $p_0$ and $\bx_1^{(j)}$ is the predicted audio for the $j$-th segment.

We introduce two extension methods below: segment-level autoregressive generation and multi-diffusion. Multi-diffusion achieves slightly better results empirically, which we use as the default method.

\textbf{Segment-level autoregressive generation.}
This algorithm emulates autoregressive generation of language models but at the segment level.
It generates one segment at a time conditioned on the information from the last $n_{ctx}$ frames of the previous segment.
Given the trajectory $\bx_{t_i}^{(j)}$ from the $j$-th segment,
we consider passing information through two routes when generating $\bx_{t_i}^{(j+1)}$: \textit{context conditioning} and \textit{trajectory regularization}.

The first route, context conditioning, is to simply update the audio context $\bc_{ctx}^{(j+1)}$ using the prediction from the previous segment.
Specifically, we set $\bc_{ctx}^{(j+1)} = [\bx_{1,-n_{ctx}:}^{(j)}; \mathbf{0}]$,
where $\bx_{1,-n_{ctx}:}^{(j)}$ denotes the last $n_{ctx}$ frames of $\bx_1^{(j)}$, and $\mathbf{0}$ is a zero matrix of shape $n_{hop} \times C$.

The second route, trajectory regularization, is to pass information through the flow trajectory $\bx_{t_i}^{(j+1)}$.
We \textit{blend} the trajectory from the previous segment with the current one at each flow step $t_i$ with a weighting function $w \in \mathbb{R}_{\ge 0}^{n_{ctx}}$,
such that the trajectory of the current segment does not diverge too much from the previous one for the overlapping segment.
Algorithmically, let $\hat{\bx}_{t_{i+1}}^{(j+1)} = \text{ODE-Solve}(u, \bx_{t_i}^{(j+1)}, \bc^{(j+1)}, t_i, t_{i+1})$ be the solution from the ODE solver at $t_{i+1}$ taking the initial condition $\bx_{t_i}^{(j+1)}$ at $t_i$ with conditioning input $\bc^{(j+1)} = \{\bc_{vid}^{(j+1)}, \bc_{ctx}^{(j+1)}, \bc_{txt}^{(j+1)}\}$ and flow model $u$.
We update the state to $\bx_{t_{i+1},0:n_{ctx}}^{(j+1)} = w \odot \hat{\bx}_{t_{i+1},0:n_{ctx}}^{(j+1)} + (1-w) \odot \bx_{t_{i+1},-n_{ctx}:}^{(j)}$ for the overlapping frames and keep other frames intact $\bx_{t_{i+1},n_{ctx}:n_{win}}^{(j+1)} = \hat{\bx}_{t_{i+1},n_{ctx}:n_{win}}^{(j+1)}$,
where $\odot$ denotes point-wise product.
In the end, we also update the audio of the overlapping frames with $\bx_{1,0:n_{ctx}}^{(j+1)}$, the consolidated prediction from the later segment.
The blending function $w$ dictates how trajectories are mixed for each frame.
To encourage a smooth transition,
we consider a linear ramping function $w_n = n / n_{ctx}$ for $n \in [n_{ctx}]$,
such that weights are higher on the previous segment for frames closer to it and vice versa.

To further boost the performance of segment-level autoregressive generation, we also explore \textit{segment-level beam search}.
At each segment, multiple candidates are generated and the resulting partial generations are ranked and pruned by a scoring model.
Those top candidates are used as prefixes to generate multiple candidates for the next segment.


\textbf{Multi-diffusion.}
Inspired by its success on leveraging a diffusion model trained on $512 \times 512$ images to generate panorama images that are 9 times wider ($512 \times 4608$) and application to video upsampling described in~\cref{sec:spatial_upsampler},
we explore multi-diffusion for audio extension, which is virtually the audio version of the panorama generation problem.
At a high level, multi-diffusion solves the ODE for each step (\eg, from $t_i$ to $t_{i+1}$) in parallel for all segments,
consolidate the prediction, and then continues with the next ODE step.
This is in contrast to segment-level autoregressive generation,
which solves the ODE for one segment completely (\ie, from $t=0$ to $t=1$, potentially with multiple steps) before solving it for the next segment.

We next describe the algorithm in detail.
To initialize, $\bx_0$ is first drawn from a prior distribution $p_0$.
At each step $t_i$, $\bx_{t_i}$ is first split into chunks $\{\bx_{t_i}^{(j)}\}_j$ based on the predefined segmentation.
$\hat{\bx}_{t_{i+1}}^{(j)}$ is computed as $\text{ODE-Solve}(u, \bx_{t_i}^{(j)}, \bc^{(j)}, t_i, t_{i+1})$ for all $j$.
Then we merge the prediction as $\bx_{t_{i+1}} = \sum_j \text{zero-pad}( m^{(j)} \odot \hat{\bx}_{t_{i+1}}^{(j)}, j)$.
The function $\text{zero-pad}(\bx^{(j)}, j)$ maps a segment of shape $n_{win} \times C$ back to a sequence of shape $N \times C$ (shape of the $\bx$),
by padding zeros based on the segment's position (\ie, the $j$-th segment spans from $n_{start}^{(j)}$ to $n_{end}^{(j)}$;
hence $n_{start}^{(j)}$ zeros are padded at the beginning and $N - n_{end}^{(j)}$ zeros at the end). 
The soft-masking function $m^{(j)} \in \mathbb{R}_{\ge> 0}^{n_{win}}$ defines how much the $j$-th segment contributes to each frame it spans.
We note that $\sum_j \text{zero-pad}( w^{(j)}, j ) = 1$, where for each frame the contribution over all segments sums to 1.
This completes one step of the flow.
The same process repeats until we reach $t_T = 1$.
Note that the weights $w$ used in segment-level autoregressive generation is applied to the overlapping frames ($n_{ctx}$ frames), 
while the soft-masking functions $\{ m^{(j)} \}_j$ are defined for all the frames each segment spans.

\cref{fig:audio_md_weights} shows two variants of soft-masking functions (zero-padded) for two segments on the right column.
The soft-masking functions $m$ are derived from window functions $\hat{m} \in \mathbb{R}_{\ge 0}^{n_{win}}$
(the left column of the same figure, also zero-padded).
Specifically, we have $\text{zero-pad}( m^{(j)}, j ) = \left( \text{zero-pad}( \hat{m}^{(j)}, j ) / \sum_{j'} \text{zero-pad}( \hat{m}^{(j')}, j' ) \right)$ given window functions $\{ \hat{m}^{(j)} \}_j$.
The uniform window function ($\hat{m}^{(j)} = \mathbf{1}$) in the top row is equivalent to the method proposed in \citet{bar2023multidiffusion}.
However, we observe that this sometimes results in abrupt transition at the boundaries (especially for smaller models),
because the prediction from two neighboring segments might still be considerably different on overlapping frames at $t=1$.
Hence the transition from taking 100\% of the first segment to 50-50 from both segments leads to discontinuity,
as shown on~\cref{fig:audio_md_weights} top row.
Using the triangular window function (\ie, Barlette window, $\hat{m}_n^{(j)} = \dfrac{2}{n_{win}-1} \left( \dfrac{n_{win}-1}{2} - \left| n - \dfrac{n_{win}-1}{2} \right| \right)$) from the bottom row results in smoother transition,
similar to what we observed in chunk-level autoregressive generation with the linear ramping function.

\begin{figure}
  \centering
  \begin{tabular}{cc}
    \includegraphics[width=0.48\textwidth]{figures/audio/md_weight_uniform_normFalse.pdf} &
    \includegraphics[width=0.48\textwidth]{figures/audio/md_weight_uniform_normTrue.pdf} \\
    \includegraphics[width=0.48\textwidth]{figures/audio/md_weight_triangle_normFalse.pdf} &
    \includegraphics[width=0.48\textwidth]{figures/audio/md_weight_triangle_normTrue.pdf}
  \end{tabular}
  \caption{\textbf{Comparing uniform weighting (top) and triangle window based weighting (bottom)} for multi-diffusion merging with an example of $l_{win}=20$ and $l_{ctx}=5$. The left column shows the unnormalized weights contributed from the first segment ([0, 15]) and the second ([10, 30]), and the right column shows the normalized weights. With triangle window, the transition in the overlapping region ([10, 15]) is much smoother.
  }
  \label{fig:audio_md_weights}
\end{figure}
